\section{Quadratic structure and localization for the new slopes}
\label{sec:analytic}

We begin the proof of Theorem~\ref{thm:local} by establishing the
quadratic detection and localization estimates for the configuration
\(y_j=x+b_jt\).  These estimates supply the input to the short-block
and frequency arguments of Section~\ref{sec:arithmetic}.  The relevant
constructions are those of \cite[Sections~3--6]{H} and
\cite[Sections~4--6]{V}.  We retain their interval projections and
stopping constructions, and prove below the coefficient-dependent
identities and estimates on which those constructions rely.

Throughout this section, the slopes are fixed.  Constants may depend on
them and on the admissibility parameters.  A rational number has height
at most \(H\) if its numerator and positive denominator in reduced form
are at most \(H\) in absolute value; matrix heights are entrywise.
We use the quantitative convention of \cite[Section~3.1]{H}: a
quasipolynomial bound in parameters \(H_1,\ldots,H_m\ge2\) is a bound
of the form
\[
 \exp\!\left(C_{q,b}\left(1+\sum_{a=1}^m\log H_a\right)^C\right).
\]
The degree \(C\) is chosen independently of \(q\), the number of
depths, and the input schedules.  A fixed number of such constructions
may be composed.  Constructions repeated a variable number of times
retain the explicit count estimates in H and V.

\subsection{Local bounds and the chart class}

For a scalar or Hilbert-valued function, write
\(\norm z_{p,I}=(|I|^{-1}\int_I\norm{z(y)}^p\,dy)^{1/p}\),
with the essential-supremum convention for \(p=\infty\).
Norms without an interval subscript use ordinary Lebesgue measure.
The elementary bound used throughout the argument is unchanged except
for its constant.
\begin{lemma}[Local absolute bound]\label{lem:local-absolute}
If \(1\le p_j\le\infty\) and \(\sum_jp_j^{-1}\le2\), then
\[
 |I|^{-2}\int_{\R^2}|w_I(x,t)|
                 \prod_{j=0}^3\norm{z_j(x+b_jt)}\,dx\,dt
 \le C_b\prod_{j=0}^3\norm{z_j}_{p_j,I}.
\]
The assertion also holds for Hilbert-valued inputs, using their norms
in the integrand.
\end{lemma}
\begin{proof}
For two distinct roles \(i,j\), the map
\((x,t)\mapsto(y_i,y_j)\) has determinant \(b_j-b_i\ne0\).
This proves the endpoint with \(p_i=p_j=1\) and the other exponents
infinite.  The support and kernel bounds give the all-infinite
endpoint, and normalized norm monotonicity gives the endpoints with
only one exponent equal to one.  Log-convexity of the positive integral
fills the stated polytope, as in \cite[Lemma~2.1]{H}.
\end{proof}

On the support of a local kernel, all four coordinates lie in an
interval of length \(r=|I|\).  Consequently
\begin{equation}
 |t|\le\frac{r}{\max_j b_j-\min_j b_j}\le\frac r3.
 \label{eq:analytic-step-support}
\end{equation}
In any pair of coordinate variables, the rescaled kernel thus has
support in a fixed box and the same type of derivative bounds.  Its
integral in either coordinate with the other fixed is zero, by the
corresponding one-coordinate marginal hypothesis.  These facts also
give \(C_N(q)\lesssim_b N\) and preserve the normalization-removal
argument of \cite[Lemma~2.4]{H}.

We recall the concrete class in which the quadratic structure is
detected.  Write \(e(u)=\exp(2\pi i u)\).  A \emph{quadratic chart atom}
is a function
\begin{equation}
 \phi(y)=e(\lambda y^2+\beta y)
       \sum_{n\in\Z^d}\chi(\theta y-n)
          e\bigl((\theta y)^TC(\theta y-n)\bigr),
 \label{eq:analytic-chart}
\end{equation}
where \(\theta\in\R^d\), \(\lambda,\beta\in\R\), \(C\in\mathbb Q^{d\times d}\),
and \(\chi\) is smooth and compactly supported.  An overall constant
of modulus one is allowed.  We use exactly the complexity convention
of \cite[Definition~4.1]{H}: for complexity \(H\ge e^2\),
\(d\le\log H\), the height of \(C\) is at most \(H\),
\(\operatorname{supp}\chi\subset[-H,H]^d\), and
\[
 \norm{\partial^\alpha\chi}_\infty
 \le H^{c_{\rm ch}(|\alpha|+1)}
             (|\alpha|+1)^{c_{\rm ch}(|\alpha|+1)}
\]
with a fixed chart constant \(c_{\rm ch}\).  The real oscillation
parameters \(\lambda,\beta,\theta\) are unrestricted.  A bounded atom
has supremum norm at most one; any normalization factor is put in its
external scalar coefficient.

The elementary operations, Weyl alternative, and chart expansion of
\cite[Lemmas~4.2--4.4]{H} concern this class itself.  Their proofs
therefore apply without a coefficient change.  We will need the precise
splitting supplied by the Weyl alternative.  In its multivariable
version replace \(\theta y\) by \(Vz\), and the ordinary quadratic
phase by \(z^T\Lambda z\), where \(\Lambda\) is real symmetric and
\(z\) has at most three coordinates.  Put
\[
 C_{\rm sy}=\tfrac12(C+C^T),\qquad
 C_{\rm sk}=\tfrac12(C-C^T),\qquad
 \operatorname{sym}M=\tfrac12(M+M^T).
\]
Let \(Q(z)\) be this multivariable chart.  Suppose \(0<\delta<1/2\)
and a Lipschitz function \(\vartheta\), supported in a fixed box
with supremum and Lipschitz norms at most \(H\), satisfies
\[
 \sup_\xi\left|\int\vartheta(z)Q(z)e(-\xi^Tz)\,dz\right|\ge\delta.
\]
The Weyl alternative then supplies a rational subspace \(U\), the
orthogonal projection \(P_f\) onto \(U\), \(P_s=1-P_f\), and the bounds
\begin{equation}
 \begin{gathered}
 (C-C^T)|_{U\times U}=0,\qquad
 V_f=P_fV,\quad V_s=P_sV,\qquad \norm{V_s}\le H',\\
 \Lambda_*=
 \Lambda+\tfrac12V^TC_{\rm sy}V
            +\operatorname{sym}(V_f^TC_{\rm sk}V_s),
 \qquad \norm{\Lambda_*}\le H'.
 \label{eq:analytic-splitting}
\end{gathered}
\end{equation}
The projection heights are at most \(H'\), and \(H'\) is
quasipolynomial in \(H,\delta^{-1}\).  The subsequent
chart expansion has the absolute coefficient and pooled Fourier-tail
bounds of \cite[Lemma~4.4]{H}.  These are the bounds used below, including
when expansions must be summed over depths.

\subsection{Separation along the configuration and local detection}

The first coefficient-dependent step moves two structured inputs into
the other two coordinates and the step variable.  This is the form of
separation needed for both detection and the later predictor estimates.

\begin{lemma}[Progression separation]\label{lem:progression-separation}
Let \(i,k,a,l\) be the four distinct roles and let \(\phi_i,\phi_k\)
be bounded chart atoms of complexity at most \(H\).  On
\(y_j=x+b_jt\), their product has an absolutely convergent expansion
\[
 \phi_i(y_i)\phi_k(y_k)
   =\sum_\nu a_\nu\psi_{a,\nu}(y_a)
                         \psi_{l,\nu}(y_l)Q_\nu(t)
\]
with bounded chart atoms as factors and the same shell bounds as
\cite[Lemma~4.5]{H}, with constants depending on \(b\).
In particular, for a quasipolynomial \(D=D_b(H)\), a dyadic Fourier
shell of size \(T\ge1\) has absolute coefficient sum at most
\(DT^{-100}\), and its factors have complexity at most \((DT)^C\).
\end{lemma}
\begin{proof}
First separate \(\phi_i\), temporarily suppressing its linear phase.
Choose the fixed rationals
\[
 \alpha_i=1,\quad\alpha_k=0,\quad
 \alpha_a=\frac{b_i-b_l}{b_l-b_a},\quad
 \alpha_l=\frac{b_a-b_i}{b_l-b_a}.
\]
They satisfy \(\sum_j\alpha_j=\sum_jb_j\alpha_j=0\).  Insert the
smooth integer partitions of unity from H at
\(\theta y_a,\theta y_l,\theta t\), and write the lifts as
\(d_j=\theta y_j-m_j\) and \(d'=\theta t-m\).  For \(j=i,l\), set
\[
 k_j=d_j-d_a-(b_j-b_a)d',\qquad k_a=0.
\]
These carries are integer vectors.  The support restrictions bound
them by the chart budget, with a fixed slope factor.  Conversely,
\(k_i,m_a,m\) determine \(m_i\) uniquely.  Each remaining specified
carry equality is imposed exactly by a smooth cutoff equal to one at
the specified integer vector and zero at all other integer vectors.

For \(S_\alpha=\sum_j\alpha_jb_j^2\), direct expansion gives
\begin{equation}
 \begin{split}
 \sum_j\alpha_j\bigl(\lambda y_j^2+(\theta y_j)^TCd_j\bigr)
  ={}&S_\alpha\bigl(\lambda t^2+(\theta t)^TCd'\bigr)\\
    &+\sum_j\alpha_j(\theta y_j)^TCk_j.
 \end{split}
 \label{eq:analytic-separation-identity}
\end{equation}
Indeed, substituting \(d_j=d_a+(b_j-b_a)d'+k_j\) cancels the
common-lift and \(x\)-terms.  Here
\(S_\alpha=(b_i-b_a)(b_i-b_l)\).  Solving for the phase in role
\(i\) leaves chart matrices \(-\alpha_aC,-\alpha_lC,S_\alpha C\)
in the three desired variables.  For fixed carries, the last sum is
linear in \(x,t\), and hence splits between \(y_a,y_l\); so does the
suppressed linear phase.

The coupled amplitude is smooth in \(d_a,d_l,d'\).  Its Fourier
expansion on a larger fixed box separates those variables, exactly as
in the proof of \cite[Lemma~4.5]{H}.  The new substitutions cost at
most \(C_b^n\) in derivatives of order \(n\), and their rational
heights grow by fixed factors.  The possible integer choices increase
by at most \(C_b^{O(d)}\).  Since \(d\le\log H\), these changes are
within the quasipolynomial chart budget.  H's integration by parts,
with derivative order exceeding the Fourier dimension by at least
\(110\), gives the stated absolute shell bound.  Separate the second
atom with \(\alpha_i=0\) and combine the factors by direct sums of
horizontal coordinates.  This completes the expansion.
\end{proof}

For one separated term, let
\[
 A_r=\sup_{\xi\in\R}
       \left|\int_\R Q(ru)\zeta(u)e(-\xi u)\,du\right|,
\]
where \(\zeta\) is smooth, nonnegative, equal to one on
\([-1/3,1/3]\), and has \(\norm\zeta_1\le1\), as in
\cite[Section~5.4]{V}.  The two-input form in the remaining roles
has norm at most \(C_bA_r\) on normalized \(L^2(I)\) inputs.
To see this, use
\[
 t=\frac{y_l-y_a}{b_l-b_a}.
\]
Fourier-expand the smooth normalized kernel in those two coordinates.
Each term is a modulated convolution with a fixed dilation of
\(Q(t)\zeta(t/r)\), so Plancherel gives the claimed bound.  If
\(A_r\ge\delta\), the Weyl alternative and chart expansion further
give the finite smooth separation in \cite[Lemma~4.6]{H}: for every
\(0<\varepsilon<1/2\), the two-coordinate kernel has a separated
approximation with uniform error at most \(\varepsilon\), whose
number of terms, atom complexities, and absolute coefficient sum
are quasipolynomial in \(H,\delta^{-1},\varepsilon^{-1}\).
The retained phases are
\(\lambda_*t^2+\gamma t\), with \(|\lambda_*|r^2\) bounded by the
stated budget, and the linear phase splits as
\[
 e(\gamma t)=e\left(-\frac{\gamma y_a}{b_l-b_a}\right)
              e\left(\frac{\gamma y_l}{b_l-b_a}\right).
\]
Thus the fixed rational-height description of the effective quadratic
and linear frequencies in that lemma is preserved.

\begin{lemma}[Local detection]\label{lem:local-detection}
Fix \(q>2\) and \(C_1\ge1\).  Suppose that the four inputs
\(R,h,u,v\), in any slot order, satisfy
\[
 \norm R_{2,I}\le C_1,\qquad \norm h_{q,I}\le1,\qquad
 \norm u_{2,I},\norm v_{2,I}\le1,
 \qquad |H_I(R,h,u,v)|>\delta,
\]
where \(0<\delta<1/2\).  There is a bounded chart atom \(\phi\)
of complexity at most \(D\) such that
\[
 \left|\frac1{|I|}\int_I R(y)\overline{\phi(y)}\,dy\right|
 \ge D^{-1},\qquad
 D\le\exp\!\left(C_{q,b}
          (1+\log C_1+\log(1/\delta))^C\right).
\]
The degree \(C\) is independent of \(q\).
\end{lemma}
\begin{proof}
We verify the progression estimate in \cite[Lemma~4.11]{H} and then
apply its detection proof.  Normalize \(I=[0,1)\).  Extend the kernel
to the same fixed circle used there.  By
\eqref{eq:analytic-step-support}, its real support fits in the
periodic \((x,t)\)-box, and all real \(y_j\) on that support already
belong to \(I\).  Hence zero-extended inputs on the circle give the
correct original integral, even when some slopes are large or negative.
The kernel's absolute Fourier coefficient sum is uniformly bounded.

The integer left inverse in Lemma~\ref{lem:algebra} distributes a
circle character with indices \((p,q')\) into slot characters with
indices \(p\mathbf1_{j=0}+q'e_j\).  To a character-weighted
progression average apply H's three Cauchy--Schwarz steps, taking
Haar differences in the directions \((-b_j,1)\).  The last cube has
three independent increments multiplied by nonzero differences
\(b_i-b_j\).  Each such integer multiplication pushes circle Haar
measure to itself.  The last integral is therefore the eighth power
of the same circle \(U^3\) norm as in H.  Character multiplication
preserves that norm, the \(L^2\) norm, and the supremum norm.  This
proves H's local generalized von Neumann estimate; its accompanying
two-\(L^1\) bound follows from Lemma~\ref{lem:local-absolute}.

The continuous inverse theorem and bounded-input decomposition in
\cite[Sections~4.5--4.6]{H} concern functions on this circle and its
chart atoms, and require no change.  In the proof of
\cite[Lemma~4.12]{H}, clip the three opposing inputs at height \(B\).
The local absolute bound gives total error at most
\(C_bC_1B^{-(q-2)/q}\).  Taking
\(B\) of order \((C_bC_1/\delta)^{q/(q-2)}\), apply that
bounded-input decomposition twice.  The error estimates use the
generalized von Neumann and two-\(L^1\) bounds just proved.
Two bounded atoms remain opposite \(R\) and an input of normalized
\(L^2\) norm at most one.  Lemma~\ref{lem:progression-separation}
and the preceding single-scale separation give a finite sum of
products of one-coordinate moments, with controlled absolute
coefficient sum.  One moment of \(R\) consequently has the asserted
size.  These are a fixed number of quasipolynomial constructions.
The clipping exponent contributes only a \(q\)-dependent multiplier
to a logarithm, so the final logarithmic degree is independent of
\(q\), as required.
\end{proof}

\subsection{The base reduction and inputs that vary with depth}

Detection is used in H and V to construct finite approximation spaces
on interval trees.  We recall their meaning before verifying the scale
estimates.  A column with label \(v\), born on an interval \(J_v\),
is an augmented function
\[
 \Phi_v(y)=(\phi_v(y),\sqrt{\epsilon_v}\,\mathbf u_v),
 \qquad y\in J_v,
\]
where \(\phi_v\) is a scalar chart atom, \(\epsilon_v>0\), and the
\(\mathbf u_v\) are distinct orthonormal private coordinates.  Only
the first, physical coordinate enters \(H_I\).  A local projection
on \(I\) is the orthogonal projection onto constant-coefficient
combinations of the available columns restricted to \(I\).  Lists
are inherited on descendants.  Accordingly the direct-sum projections
on successively finer interval partitions have nested ranges.  The
private-coordinate, compression, and height-bin estimates of
\cite[Section~3]{H} depend only on these definitions and the column
budgets, and apply unchanged.

There are two coefficient-dependent scale estimates in the base
construction.  For two fixed chart atoms, the separation above reduces
the other two inputs to a kernel of the form
\(w_I(x,t)Q(t)\), with fixed bounded one-coordinate multipliers.
The proof of \cite[Lemma~5.1 and Corollary~5.2]{H} expands \(Q\),
groups nearby linear frequencies, and uses the zero marginals of
\(w_I\).  In coordinates \((y_a,y_l)\), its double primitive has
compact support in an interior rectangle.  Fourier expansion followed
by one derivative in each coordinate expresses it as a sum of tensor
products of mean-zero smooth bumps.  In role \(p\in\{a,l\}\), the
resulting scalar rows have fixed bounded carrier functions
\[
 m_{\nu,p}(y)=e(\nu y)\psi_p(y),
 \qquad \nu=\pm\frac{\gamma}{b_l-b_a},
\]
where \(\psi_p\) is the fixed separated multiplier in that role.
These are the only coefficient replacements in the row estimate.
They alter frequency separations and rational heights by fixed
factors.  H's height-bin argument, ordinary Bessel estimate, and
compression step therefore give its absolute scale sum and finite
predictor list with the same quantitative orders, including the
pooled-tail bound needed for summation over all depths.  The finer
partition comparison in \cite[Lemma~5.3]{H} uses the resulting
one-coordinate predictor moments, so its projection-increment proof
applies with the same replacement.

For the second estimate, fix three bounded chart columns and take
the transpose in role \(l\) at \(y+h\).  Its other coordinates are
\[
 y+h+(b_j-b_l)t\qquad(j\ne l).
\]
After squaring the normalized \(L^2\) norm, the six atoms form a
multivariable chart in \((h,t,t')\).  At a fixed path point \(y\),
all parameters before multiplication by \(r\) or \(r^2\) remain
fixed as the depth changes.  H's Weyl alternative, chart expansion,
and height-bin argument therefore apply.  Away from the bounded
number of intermediate scales, the bounded effective quadratic
parameters are tiny and each retained \(t\)-frequency is tiny or
large.  Smoothness handles the large frequencies.  Removing the tiny
oscillations leaves exactly
\[
 \int_\R w_I(y+h-b_lt,t)\,dt=0.
\]
This proves the active-scale count and the small-norm scale sum in
\cite[Lemma~5.7]{H}, with constants depending on \(b\).

For scheduled inputs we use \cite[Section~4]{V}.  Its row estimates
concern Hilbert-space operators on refining interval partitions,
evaluated on inputs at prescribed nondecreasing read depths.  Their
hypotheses are support on the input partition, representation of the
adjoint ranges on the observation partition by the inherited columns,
the stated order of successive partitions, and an ordinary Bessel
bound before evaluation on the varying inputs.  None involves the
maps \(y_j\).  Restriction to a new spatial root preserves the
pointwise block bound in Theorem~\ref{thm:local}.  The original
\(W_j\) and original depth schedules are retained; subsequences may
have empty intersections with some blocks, as allowed in V's row
argument.  Its nearest-endpoint extension is retained as well.

The smooth rows in \cite[Lemma~4.6]{V} require the ordinary Bessel
bound just verified, with the fixed carriers above.  Thus the
scheduled scale-sum and predictor results
\cite[Lemma~5.3 and Corollary~5.4]{V}, and its delayed replacement
Lemma~5.5, hold with the new coefficients.  The remaining steps of
its Section~5 use the same local detection, fixed-cell fits,
projection energies, and staggered pair comparisons.  A finer
projection in a comparison is still applied to the input at the
prescribed read depth, and its row estimate has no factor depending
on \(W_j\).

We now describe the output of the base construction.  For the
fixed-input argument, embed the scalar input as
\(F_{j,s}=(f_j,0)\), independent of \(s\), on a subtree satisfying
\(\norm{f_j}_{q,I}\le1\) on every cell, as in
\cite[Proposition~2.3]{H}.  For the scheduled argument, put
\(F_{j,s}=(z_{j,s},0)\) with the inputs of Theorem~\ref{thm:local}.
The adaptive fits and two-earliest-increments expansion of
\cite[Lemmas~5.4--5.5]{H}, or \cite[Sections~5.1--5.2]{V}, express
the original form, at any finite accuracy cutoff, in tuples with
two differences of fitted projections and two original-input
residuals, together with a remainder whose absolute length-weighted
sum tends to zero on the fixed finite tree.  Each tuple has an
earlier and a later accuracy index, and a choice of the four roles.  Their
exact staggered ordering is retained from those results.
High-position truncation, using the predictor estimates above, is
\cite[Proposition~5.6]{H} or \cite[Sections~5.5--5.6]{V}.  It assigns
each original summation cell to a segment root \(R\) with inherited
projection families; write \(\mathcal I_R\) for the cells assigned
to \(R\).  Auxiliary projection cells are not additional summands.

Fix one tuple and let \(k\) be its later accuracy logarithm.
The two positions called low in this construction are denoted
\(i,j\), and the other two by \(a,l\).  On a cell of depth \(s\),
the resulting arguments satisfy
\begin{equation}
 \begin{gathered}
 Y_{p,I}=B_{p,I}+U_{p,I},\qquad U_{i,I}=U_{j,I}=0,\\
 B_{p,I}=\sum_{h=1}^{m_p}\varepsilon_{p,h}Q^{p,h}_{s}F_{p,s},
 \qquad m_p\le2,\quad \varepsilon_{p,h}\in\{-1,1\},\\
 U_{p,I}=(P^{\rm new}_{p,s+K_0}-P^{\rm new}_{p,s})F_{p,s}
                         \quad(p=a,l).
 \end{gathered}
 \label{eq:analytic-basic-arguments}
\end{equation}
All expressions are restricted to \(I\); a projection indexed by
a finer depth is the direct sum on that finer partition.  The
\(Q^{p,h}\) and \(P^{\rm new}_p\) are the inherited basic families
just constructed, and \(K_0\) is their positive integer delay.
Their column sizes, inverse private weights, scalar complexities,
coefficient sums, and path counts have a joint budget
\(e^k\le D_{\rm base}\le\exp(C_{q,b}k^C)\).
On every segment they give
\begin{equation}
 \begin{gathered}
 \norm{B_{p,I}}_{2,I},\ \norm{U_{p,I}}_{2,I},\
 \norm{Y_{p,I}}_{2,I}\le C_{q,b},\qquad
 \Lambda_k\le C_{q,b}k^C,\\
 \sum_{I\in\mathcal I_R}|I|\norm{U_{p,I}}_{2,I}^2
                  \le\Lambda_k^2|R|\quad(p=a,l).
 \end{gathered}
 \label{eq:analytic-basic-energy}
\end{equation}
One may take \(\Lambda_k\) to dominate
\(C_{q,b}(1+K_0)\log(2D_{\rm base})\).  For the fixed-input
argument the smaller bound in H also suffices.

For noninitial accuracy levels, every pair of full \(Y\) arguments
has form norm at most \(\eta=\exp(-k^\theta)\), and
\(|H_I(Y_I)|\le\eta\).  Here the pair norm is the supremum against
normalized \(L^2(I)\) inputs in the complementary two positions;
\(0<\theta<1\) is fixed before \(q\).  Define \(\mathcal A_R\)
to consist of those cells for which some available triple of basic
scalar columns has transpose norm greater than
\(D_{\rm base}^{-C_{\rm act}}\), with the fixed sufficiently large
exponent in H or V.  The active-scale estimate gives
\(\sum_{I\in\mathcal A_R}\mathbf1_I\le D_{\rm base}^C\).

There are three remaining patterns: \(z^\varnothing=B\),
\(z^a\) replaces only role \(a\) of \(B\) by \(U_a\), and
\(z^l\) replaces only role \(l\) by \(U_l\).  Put
\begin{equation}
 \mathfrak M_R(z)=\sum_{I\in\mathcal A_R}|I|
                         \min\{\eta,|H_I(z_I)|\}.
 \label{eq:analytic-base-target}
\end{equation}
The output of \cite[Proposition~5.8]{H} and
\cite[Proposition~5.6]{V} is
\begin{equation}
 \sum_{I\in\mathcal I_R}|I|\,|H_I(Y_I)|
 \le C_{q,b}\eta\Lambda_k^2|R|+e^{-k}|R|
              +\sum_{z\in\{z^\varnothing,z^a,z^l\}}\mathfrak M_R(z).
 \label{eq:analytic-active-reduction}
\end{equation}
Indeed the term with both \(U\) factors uses the full complementary
pair bound and \eqref{eq:analytic-basic-energy}.  The inactive
triple scale sum bounds the omitted cells.  On the others, apply
\(\min(\eta,\sum x_m)\le\sum\min(\eta,x_m)\) to the full
tuple before separating its patterns.  No individual pattern is
asserted to have the full-pair smallness.

For each fixed tuple the segment roots have total length bounded
by a fixed multiple of the original root length, and the sum of
the high-truncation replacement errors is at most \(C_{q,b}e^{-k}\)
times that length.  These are the segment-packing assertions in
the cited base constructions.  The finitely many initial accuracy
levels use the active path bound and local sizes, without the
factor \(\eta\), and hence have depth-independent bounds.  All
these estimates are valid anew on a normalized subroot with its
inherited lists and original input schedule.  We have therefore
reduced both input classes to estimating
\eqref{eq:analytic-base-target} uniformly over these concrete base
patterns.

\subsection{Curvature tags and separation of distinct tags}

We now localize the base arguments according to their quadratic
structure.  Recall the nonzero integers \(c_j\) of
Lemma~\ref{lem:algebra}.  A column in role \(i\) has
\emph{normalized curvature tag} \(v=(\theta_v,C_v,\lambda_v)\)
when its quadratic chart phase is written as
\[
 c_i\bigl(\lambda_v y^2+(\theta_vy)^TC_v(\theta_vy-n)\bigr).
\]
Its linear modulation is not part of its tag.  This is H's
normalization with \((-1,3,-3,1)\) replaced by \(c\).  Division by
the fixed nonzero \(c_i\) preserves the chart budgets.

At length \(r\), the difference data of two tags \(v,v'\) are
\begin{equation}
 V'=r\binom{\theta_v}{\theta_{v'}},\qquad
 C'=\operatorname{diag}(C_v,-C_{v'}),\qquad
 \Lambda'=r^2(\lambda_v-\lambda_{v'}).
 \label{eq:analytic-tag-difference}
\end{equation}
The ambient graph in \cite[Section~6.1]{H} joins two tags when these
data admit \eqref{eq:analytic-splitting} with bounds at most the
chosen edge height \(H_{\rm ed}\); equal tags have distance zero.
This graph includes all hypothetical tags within the basic budget.
Its distances depend on tags and depth, and do not increase when
the interval length is halved.

\begin{lemma}[Separation of distinct tags]\label{lem:tag-separation}
For single-column complexities at most \(\exp(C_{q,b}k^C)\) and
\(\log(1/\tau)\le C_{q,b}k^C\), the edge height can be chosen with
\(\log H_{\rm ed}\le C_{q,b}k^C\) so that tags at graph distance
greater than one have both of the following properties.  Their
same-slot normalized physical inner product is at most \(\tau\)
in absolute value.  In distinct slots, their form norm against two
arbitrary normalized \(L^2(I)\) spectators is at most \(\tau\).
\end{lemma}
\begin{proof}
For the first assertion use the Weyl alternative on the conjugate
product, after the smoothing of the interval indicator in
\cite[Lemma~6.1]{H}.  Its fixed nonzero factor \(c_i\) changes
the allowed rational heights by a fixed amount, and multiplying
back gives the splitting for the tag difference.

For distinct source slots \(i,j\), let the spectators be at
\(a,l\), and write the centered source coordinates in terms of
the normalized spectator coordinates \(z,w\).  Their linear
parts are \(r(\xi_pz+\zeta_pw)\), where
\[
 \xi_p=\frac{b_l-b_p}{b_l-b_a},\qquad
 \zeta_p=\frac{b_p-b_a}{b_l-b_a}\qquad(p=i,j).
\]
These are nonzero fixed rationals.  The cancellation of polynomials
of degree at most two gives
\begin{equation}
 c_i\xi_i\zeta_i=-c_j\xi_j\zeta_j=:\kappa\ne0.
 \label{eq:analytic-mixed-curvature}
\end{equation}
Use the difference data \eqref{eq:analytic-tag-difference}.  For
a small return \(V'h=m+d_h\), choose an integer period \(\Pi\)
and require \(m\in\Pi\Z^{d_v+d_{v'}}\), where \(d_v,d_{v'}\)
are the horizontal dimensions of the two tags.  Choose \(\Pi\)
sufficiently divisible that \(\xi_i\Pi,\xi_j\Pi\) are integers and
\(\Pi C_v,\Pi C_{v'},c_i\xi_i\Pi C_v,c_j\xi_j\Pi C_{v'}\)
have integral entries.
This clears the lattice-index and phase denominators together.
Translate \(z\) by \(h\), and the source lift indices by
\(\xi_im_v\) and \(\xi_jm_{v'}\), respectively.  The phase-increment
calculation in the proof of \cite[Lemma~4.3]{H}, together with
\eqref{eq:analytic-mixed-curvature}, gives the modulation in \(w\)
\[
 \kappa\Omega(h,m),\qquad
 \Omega(h,m)=(2\Lambda'+(V')^TC'V')h
                      -(V')^T(C'-(C')^T)m.
\]
Indeed each source contribution has the multiplier
\(c_p\xi_p\zeta_p\); the two contributions therefore combine
into these difference data.  The remaining linear terms in \(z\)
are absorbed in the first spectator.  The amplitude and kernel
errors are bounded by a quasipolynomial times
\(|h|+\norm{d_h}\), as in H.

Call the two spectators \(f,g\), and clip them at a sufficiently large inverse polynomial
in \(\tau\).  Their discarded \(L^1\) tails are small by their
\(L^2\) bounds and Lemma~\ref{lem:local-absolute}.  If the form is
larger than \(\tau\), applying the return calculation to sufficiently
many multiples of \((h,m)\) gives large correlations of one fixed
bounded kernel with tests \(f_\ell(z)g_\ell(w)\), where the
\(f_\ell\) have bounded \(L^2\) norms and
\[
 g_\ell(w)=g(w)e(\ell\kappa\Omega(h,m)w)
\]
up to a common bounded multiplier.  The small return threshold is
chosen before these multiples so that all translation errors remain
small.  No regularity of the first spectator is required: it is
translated by a change of variables.  If \(|\Omega(h,m)|\) exceeds
a fixed constant, the second tests are Bessel on the fixed box,
with bound controlled by the clipping height.  Cauchy--Schwarz in
\(z\) then contradicts the number of large correlations.  Thus
\(\Omega\) is bounded on every sufficiently small return.

The concluding return-to-splitting argument in the proof of
\cite[Lemma~4.3]{H} uses precisely this bound: its alternating
return identity yields an isotropic lattice space, successive
minima bound the rational fast projection and slow part, and
coordinate-direction returns bound the effective quadratic form.
The enlarged denominator-clearing period changes only its
quasipolynomial constants.  It supplies
\eqref{eq:analytic-splitting}, contradicting the assumed graph
distance.  This is the distinct-slot argument of
\cite[Lemma~6.1]{H} for the present coefficients.
\end{proof}

\subsection{Same-tag tests and the localization output}

The complementary fact is that a column of one tag can detect a
residual in another role using a template of that same tag.  This
keeps the approximation spaces finite without repeatedly increasing
their nonlinear complexity.  A pure source column means one scalar
column label, rather than a linear combination of labels.

\begin{lemma}[Uniform same-tag test]\label{lem:same-tag-test}
Fix an original tag \(v\) and its original source columns within
the complexity budget \(\exp(C_{q,b}k^C)\).
There are fixed template families in
all four slots, consisting of bounded columns of tag \(v\) with
arbitrary external linear modulations.  Suppose a pure source
column from the original or template family in role \(i\) has
form norm greater than \(t_0>0\) against a physical residual in \(L^2(I)\) in
role \(j\) and two normalized \(L^2(I)\) spectators.  The residual
then has a normalized moment of magnitude at least \(t_0/H_0\)
against a template in role \(j\), where
\(H_0\le\exp(C_{q,b}k^C)\).  This bound is independent of the
external linear modulations and of the number of earlier template
insertions.
\end{lemma}
\begin{proof}
Use the nonnegative partition function \(\rho\) and template
amplitudes \(\rho(d)e(\ell^Td/P_0)\) of
\cite[Lemma~6.6]{H}, with \(\ell\) ranging over residue vectors
modulo a sufficiently large integer period \(P_0\).  For the four
lifts \(d_p=\theta_vy_p-n_p\), set
\begin{equation}
 d_\diamond=\sum_pe_pd_p,\qquad
 k_p=d_p-d_0-b_pd_\diamond.
 \label{eq:analytic-common-carries}
\end{equation}
The B\'ezout identities show that \(d_\diamond\) is a lift of
\(\theta_vt\) modulo the integer lattice and every \(k_p\) is
integral.  All are bounded on the branch boxes.  The common
quadratic phases cancel by Lemma~\ref{lem:algebra}; for fixed
carries the remaining phase is
\[
 \sum_pc_p(\theta_vy_p)^TC_vk_p,
\]
which is linear in the coordinates.  Smooth equality cutoffs in
the four lifts impose \eqref{eq:analytic-common-carries} exactly,
because their arguments are integer valued on the configuration.
Their Fourier expansion on a sufficiently large box therefore
produces the same template amplitudes.  Increasing the period
and the possible carry count costs only the allowed chart budget.

As in H, write each Fourier index as \(\ell=P_0q'+r'\).
The factor \(e((q')^Td_p)\) is an external linear modulation,
so only \(r'\) enters the template amplitude.  This gives the
same bound even when the source was itself a previously inserted
template.  For an original source, the remaining smooth horizontal
amplitude, after phase cancellation, is periodized and expanded
in ordinary torus characters.  Its remaining dependence on the
source coordinate is a linear modulation, which is absorbed into
the spectator tests.  All these absolute coefficient sums are
bounded by the same quasipolynomial budget.
After separating the smooth kernel, express the target
coordinate as \(y_j=\alpha y_a+\beta y_l\), with fixed nonzero
rational \(\alpha,\beta\).  Fourier-expand the demodulated
residual on a fixed larger interval.  If \(M\) is the supremum
of its normalized template moments, one separated term is bounded
by
\[
 M\sum_m|A_mB_m|
 \le M\left(\sum_m|A_m|^2\right)^{1/2}
       \left(\sum_m|B_m|^2\right)^{1/2}
 \le C_bM.
\]
The last inequality is Bessel's inequality for spectator
frequencies with nonzero spacings proportional to
\(\alpha,\beta\).  Summing the absolute Fourier coefficients
gives \(H_0M\), and proves the assertion.
\end{proof}

We record how these two tests enter the remainder of the imported
construction.  This also fixes the meaning of its output used in
Section~\ref{sec:arithmetic}.  Put \(P=\lceil\eta^{-\alpha}\rceil\),
where \(\alpha>0\) will be chosen in the closing argument.  The
vertices are occurrences of the original columns and are ordered by
admission.  Added main and helper columns attach to those vertices
and have their normalized tags; they create no new graph vertices.
H's short-path, color, and historical-mask constructions
\cite[Lemmas~6.2--6.5]{H} form groups of these vertices and assign
current groups to each summation cell.  An actual use is such an
occurrence of a current group at a cell.  A mask is
the subset of available vertices allowed in a specified local
projection.  A pad enlarges that subset by the prescribed graph
radius; a positive smeared filter is the average of \(P\) nested
pad projections.  A history is the ordered product of such filters
or their complements used in the binary localization.  The number
\(\mathcal L\) of colors is polynomial in \(k\), and these main
histories have \(O(\log\mathcal L)\) factors.
Each group records a distinguished tag \(z\), the primary winner
in H's color construction.  Every vertex admitted to its masks
has, already at admission, a graph path to \(z\) of length
polynomial in \(P,k\).  The short-path lemma therefore gives the
endpoint splitting \eqref{eq:analytic-splitting} relative to \(z\),
with logarithmic height polynomial in \(P,k\).  This is the
splitting used for the first matching of a column in
Section~\ref{sec:arithmetic}.

We use the precise masks and historical-use convention of those
lemmas.  Its relevant properties are as follows.  Masks are
inherited on descendants; at an actual use, enlarging a pad by one
includes every neighbor of the preceding mask.  An endpoint is usable when its
mask is nonempty and its cell has a downstream actual use of its
group.  Distinct groups of one color have tags separated by graph
distance greater than one on overlapping usable endpoints.
The endpoint bound along short paths follows from the rational
splitting \eqref{eq:analytic-splitting}.  The usable-space Bessel
bound follows from same-slot separation in
Lemma~\ref{lem:tag-separation}.  These proofs involve only rational
linear algebra, inherited interval projections, and the indicated
splitting and separation bounds; thus their path, rank, and pad
counts keep the quantitative orders in H.

For completeness, the closure condition imposed on these spaces
has an exact form.  At rank \(r\), let \(n_{r-1}\) bound the
number of labels through the preceding rank along a path, let
\(D=D_{\rm base}\), and let \(\tau_3=D^{-C_\tau}\).  For a
basic argument \(a\), each prescribed history \(\mathcal H\),
target mask \(\mathcal M\), and pure lower-rank source whose
tag belongs to that mask, the form norm against the residual
\((1-P_{r,\mathcal M})\mathcal Ha\) and two normalized
\(L^2\) spectators is required to be at most
\begin{equation}
 \frac{\tau_3}{D\sqrt{n_{r-1}}}\norm a_{2,I}.
 \label{eq:analytic-closure}
\end{equation}
On failure, Lemma~\ref{lem:same-tag-test} supplies a new column
and hence a projection increment.  The products-of-prefixes
estimate of H Section~3 bounds the sum of those increments,
exactly as in its Lemma~6.7.  Its rank recurrence and main/helper
catalog counts are therefore unchanged.  This count depends on the
names of masks and histories, including anticipated helper bases,
but not on the numerical edge height or on later mode precision.
As in \cite[Section~6.4]{H} and \cite[Section~6.1]{V}, first fix
these counts, then the mismatch tolerance, and finally the edge
height needed in Lemma~\ref{lem:tag-separation}.
The lower-rank count in
\eqref{eq:analytic-closure} is retained, so no full source count
is lost when these errors are summed.

It follows that the binary localization and outgoing-filter
identities of \cite[Proposition~6.8 and Lemma~6.9]{H} hold for
the present forms: they use \eqref{eq:analytic-closure}, distinct-tag
form separation, the local absolute bound, and the same nested
projection algebra.  In the color construction of
\cite[Lemma~6.3]{H}, the first failure of its marking conditions
may remove a subtree.  Increasing the fixed thresholds makes the
sum of these roots' lengths an arbitrarily small fixed fraction
of \(|R|\).  Section~\ref{sec:completion} restarts the original
summands on those roots.  After these marking-failure cuts, let
\(\mathcal A_R^{\rm main}\) denote the remaining active cells.
For each base pattern and \(I\in\mathcal A_R^{\rm main}\) of depth
\(s\), the localization is
\begin{equation}
 \begin{gathered}
 H_I(z_I)=\sum_\lambda\ \sum_{G\ {\rm current}}
       H_I(Z^\lambda_{0,G,s},\ldots,Z^\lambda_{3,G,s})+e_I,\\
 Z^\lambda_{p,G,s}
       =\overline L^\lambda_{p,G,s}\mathcal B^\lambda_{p,s}z_{p,I},
 \qquad
 \sum_{I\in\mathcal A_R^{\rm main}}|I||e_I|
           \le C_{q,b}e^{-k}|R|.
 \end{gathered}
 \label{eq:analytic-localization}
\end{equation}
Here \(\lambda\) ranges over polynomially many fixed binary search
leaves.  The inner history \(\mathcal B^\lambda_{p,s}\) is independent
of \(G\), and \(\overline L^\lambda_{p,G,s}\) is the outer ball
average.  The minimum in \eqref{eq:analytic-base-target} is
distributed across the fixed leaves, while the group sum remains
inside the original cell contribution.  The filters act on the
whole basic arguments in this identity.  We next record their
endpoint bounds, then the algebra that moves the complementary
filters onto a side operand.

For varying inputs, this localization acts at each cell on its
one actual tuple at that cell's original depth.  The insertion
tests freeze that operand while the column list grows.  The
products-of-prefixes estimate therefore remains a fixed-vector
estimate, and gives the same closure count; this is the argument
of \cite[Section~6.1]{V}.  After the operator families have been
fixed, its Lemma~6.1 and Corollary~6.2 control their changes between
endpoint depths.  More explicitly, for one slot, search leaf,
and color, write the full localized endpoint output as
\(Z_{G,s}=\overline L_{G,s}\mathcal B_s B_s\), where
\(\overline L\) is the outer positive ball average and
\(\mathcal B\) its group-independent inner history.  If depth
intervals \([r,t]\) have edge-overlap multiplicity at most \(a\),
and \(E_{G,r,t}\) is a union of \(t\)-cells with downstream actual
uses certifying the indicated group, those results give
\begin{equation}
 \sum_{G,[r,t]}
  \norm{\mathbf1_{E_{G,r,t}}(Z_{G,t}-Z_{G,r})}_2^2
 \le aG_0|R|,
 \qquad G_0=\exp\!\bigl(C_{q,b}(1+\log(2k))^C\bigr).
 \label{eq:analytic-endpoint-variation}
\end{equation}
The shared histories without an outer ball satisfy the corresponding
estimate counted once across groups.  Endpoints before the attempt
root are zero, and all other endpoints use inherited continuations
and local refits.  The proof first establishes the ordinary Bessel
bounds for operator changes, using usable endpoint ranges and
compression, and then evaluates the rows at their prescribed
depth inputs.  Input changes use only the original block variation
bound.  Thus the constants in
\eqref{eq:analytic-endpoint-variation} are independent of the
partitions and \(W_j\).

\subsection{Transferring the complementary filters}

The side estimate will encounter two distinguished operands while
the other two retain their whole localized histories.  We state
the projection identity for this situation explicitly.  Fix an
actual group use at a cell \(I\), one search leaf, and four distinct
roles \(a,a',c,d\).  Let \(S_a,S_{a'}\) be fixed Hilbert-valued
\(L^2(I)\) operands in the first two roles.  They are called side
operands here because they are retained while the complementary
histories are simplified; their particular construction and
cross-depth bounds are given in Section~\ref{sec:arithmetic}.
The other two roles have the whole localized arguments with basic
inputs \(B_c,B_d\) from \eqref{eq:analytic-localization}.
Choose the side operands before closing the helper catalog.
The helper version of \eqref{eq:analytic-closure} is imposed with
these fixed states as bases and lower-rank main columns as sources,
as in \cite[Lemma~6.7]{H}.

Expand each complementary factor \(1-\overline Q\) as the
difference of the identity and a positive average.  There are only
\(2^{O(\log\mathcal L)}\) resulting products.  In one such product,
remove the largest-rank outermost positive filter in roles \(c,d\).
The positive main ranks are distinct in the source construction,
so the remaining source in the other complementary role has
strictly lower rank.  Suppose the filter being removed is on
role \(c\).  Let \(Q_{c,p}\), \(1\le p\le P\), be its individual
pad projections, and \(Q_{a,p}\) the helper-rank projections in
role \(a\) with the same masks and pads.  For \(j=a,c\), set
\begin{equation}
 \begin{gathered}
 Q_{j,0}=0,\qquad Q_{j,P+1}=1,\qquad
 A_j^p=Q_{j,p}-Q_{j,p-1}\quad(1\le p\le P+1),\\
 \overline Q_j=\frac1P\sum_{u=1}^P Q_{j,u}
       =\sum_{p=1}^{P+1}\frac{P-p+1}{P}A_j^p .
 \end{gathered}
 \label{eq:analytic-bands}
\end{equation}
The \(A_j^p\) are mutually orthogonal projections and sum to the
identity.  Their final band is included, with weight zero in
the positive average.

Let \(x_c\) be the inner history on \(B_c\) after deleting the
outgoing filter, and let \(x_d\) be the entire remaining history
on \(B_d\).  Write \(x_a=\mathcal D_aS_a\), where \(\mathcal D_a\)
is the product of helper averages already transferred, and
\(x_{a'}=S_{a'}\).  In the next displays arguments are written in
role order \(a,a',c,d\); each still occupies its original slot in
\(H_I\).  Multilinearity and \eqref{eq:analytic-bands} give
\begin{equation}
 \begin{split}
 &H_I(x_a,x_{a'},\overline Q_cx_c,x_d)
       -H_I(\overline Q_ax_a,x_{a'},x_c,x_d)\\
 &\qquad=\sum_{p,q=1}^{P+1}\frac{p-q}{P}
                  H_I(A_a^px_a,x_{a'},A_c^qx_c,x_d).
 \end{split}
 \label{eq:analytic-band-transfer}
\end{equation}
Indeed the coefficient of the pair of bands is the weight of
band \(q\) in role \(c\) minus the weight of band \(p\) in
role \(a\).  This proves H's Equation~(15), including its
terminal bands, for the present form without a coefficient change.

The estimate accompanying \cite[Lemma~6.9]{H} permits only
\(|p-q|\le2\) to be retained: the other terms at this cell and
group have total absolute value at most
\[
 D_{\rm base}^{-C_3}
       \norm{S_a}_{2,I}\norm{S_{a'}}_{2,I}
\]
for any prescribed fixed \(C_3\), after the numerical accuracies
are chosen.  Its hypotheses are precisely the closure and
separation estimates proved above.  If \(p>q+2\), expand the
lower \(c\)-band into pure main sources and use the helper
closure test on the higher \(a\)-band, written as a difference
of residual projections.  If \(q>p+2\), expand the lower-rank
source in role \(d\).  A source tag within one edge of the
lower helper mask belongs, by neighbor closure, to the next
outgoing mask; the main closure test then controls the higher
\(c\)-band.  For a tag farther away, expand the helper primal
and use Lemma~\ref{lem:tag-separation}.  The coefficient counts
are canceled by the lower-rank denominator in
\eqref{eq:analytic-closure}, except for fixed count powers
already included when the mismatch tolerance was chosen.
The \(O(P^2)\) band pairs and polynomially many transfer
positions fit those accuracy choices because \(\log P=O(k)\).

For each fixed offset \(h=q-p\in\{-2,-1,0,1,2\}\), retain the
coupled sum
\begin{equation}
 \mathcal B_{I,h}
    =-\frac hP
       \sum_{\substack{1\le p\le P+1\\1\le p+h\le P+1}}
          H_I(A_a^px_a,x_{a'},A_c^{p+h}x_c,x_d).
 \label{eq:analytic-retained-bands}
\end{equation}
The sum over \(p\) stays inside the cell contribution.  When the
current group is also summed, its sum stays there as well.
This is the expression to which the short-block estimate is
applied; taking absolute values of its individual bands would
lose the matched-band summation used in that estimate.

Forward only the transferred main term in
\eqref{eq:analytic-band-transfer} to the next step.  Removing
the remaining positive filters in decreasing rank ends at
\begin{equation}
 H_I(\mathcal D_aS_a,S_{a'},B_c,B_d),
 \qquad \mathcal D_a
   \text{ a product of }O(\log\mathcal L)\text{ helper averages}.
 \label{eq:analytic-transfer-main}
\end{equation}
Restoring \(B_c,B_d\) to the full \(Y_c,Y_d\) uses their pair
bound from the base reduction.  Each correction has an additional
basic \(U\) factor and is therefore a three-side term.

For the retained defects, the endpoint convention matters.
If the outgoing filter is a ball filter, split
\(A_c^{P+1}=1-Q_{c,P}\) into the shared identity and its positive
ball endpoint.  The other complementary band endpoints have
usable ball ranges.  If the outgoing filter is cumulative,
its outer ball has already been removed and the entire
complementary band expression is shared across groups.
Role \(d\) likewise either has one positive outer ball filter
or is shared.  In the endpoint energy estimates, a shared
identity is evaluated on the union of required cells and counted
once; the original form still retains every actual use.
These are exactly the endpoint
classes controlled by \eqref{eq:analytic-endpoint-variation}
and \cite[Corollary~6.2]{V}; they give the side estimate in
Section~\ref{sec:arithmetic} its coupled band operands.

We have obtained the base minimum, the localized group expressions,
their outgoing-filter identity, and the endpoint estimates with the
same main-column budgets as H and V.  The remaining work is to
estimate the coupled near bands and the remaining localized quadratic
contributions.  Section~\ref{sec:arithmetic}
establishes the coefficient-dependent comparisons used for those
two estimates.
